% Cote 119: offprints, reports, courses, CV and letters (typed and printed part): archive cote 119, site page 262 (PDF page 263).
% typescript. Draft compared: none.
% AI-transcribed by Claude Opus 5.5 (claude-opus-5-5, adaptive thinking, effort high), Message Batches API, batch r119_full (msgbatch_01DwZNm5Y7YsKBWyAN4sFs8w), 2026-11-13; not yet checked by eye.
\scanpage{119/262}
\mspage{9}

conditions favorables à son épanouissement qui ont jusqu'à présent fait défaut.

En plus des maîtres eux-mêmes, les étudiants ont été évidemment les
principales victimes de la survie d'un système de formation et de recrutement
des maîtres basé presque exclusivement sur le critère de l'aptitude à la
recherche. Il est urgent de mettre à profit dans l'immédiat la force vive
libérée aujourd'hui par le mouvement de revendication étudiantes, et les
prises de conscience qu'elles ont précipitées chez les enseignants, pour
mettre au point des mécanismes assez souples pour distinguer et coordonner
tout à la fois les deux principales vocations de l'Université, et capables
de s'adapter au fur et à mesure aux besoins, changeants d'une période à l'autre
et d'une Faculté à l'autre. Ces questions sont l'objet dès à présent de
discussions entre étudiants et enseignants (assistants, maîtres-assistants,
maîtres de conférences et professeurs), tout au moins au département de
Mathématiques de la Faculté des Sciences d'Orsay, et sans doute un peu partout
dans les universités en effervescence créatrice. Il faut que nous, les
principaux concernés, étudiants et maîtres enseignants et chercheurs, fassions
tout notre possible pour que des solutions d'ensemble soient dégagées dans
le plus proche avenir, et puissent être mises en application dès la rentrée
prochaine. Je passerai en revue plus bas diverses modifications en profondeur
des structures et des habitudes de pensée dont la nécessité me semble dès
à présent acquise. Certaines parmi les plus importantes découlent logiquement
de l'é\xout{\unsure{t}}tat de faits que je viens d'essayer d'esquisser. Cet état de faits
est bien familier sans doute à ceux de mes collègues ayant assez d'expérience
de l'enseignement ou de la recherche, mais moins familier déjà, semble-t-il,
aux assistants ou aux étudiants, et a fortiori aux personnes étrangères à
l'Université. Mon but dans la première partie de cet exposé a été de mettre
en évidence le plus clairement possible un des vices de base du système actuel,
et d'inciter tous ceux qui sont concernés à réfléchir à cette question pour
contribuer à l'élaboration de solutions novatrices. Les incertitudes du moment
présent sont plus lourdes de promesses que toutes les certitudes sur lesquelles
se reposait dame Université jusqu'à il y a peu. C'est de nous tous qui sommes
concernés, étudiants, enseignants et chercheurs, non des princes qui nous
gouvernent, qu'il dépend que ces promesses débouchent sur cette rénovation
de l'enseignement dont personne (sujets ni princes) ne songe plus à contester
la nécessité.

\margin{262}
